\section{Architektur und Technologiestack}
\label{chap:architektur}

Drei Eigenschaften des MVP prägen die Architektur. Die App verarbeitet Hautfotos und damit Gesundheitsdaten im Sinne von Art.~9 DSGVO~\cite{dsgvo}. Der Hauptanwendungsfall hängt von Kamera und Bildanalyse ab. Und das MVP sieht kein Nutzerkonto vor. Daraus folgt die Grundentscheidung: Die App arbeitet vollständig auf dem Gerät und kommt ohne eigenes Backend aus.

\subsection{Komponenten}

Aus den Anwendungsfällen ergeben sich acht Komponenten: Benutzeroberfläche, Kameraaufnahme, Qualitätsprüfung, Eignungsprüfung, KI-Einstufung, Datenhaltung, Körperübersicht und Erinnerungen. Welche Anwendungsfälle jede davon abdeckt, zeigt \tabref{stack}.

Die App ist in drei Schichten gegliedert:
\begin{itemize}
\item Die \emph{UI-Schicht} enthält Screens und ViewModels.
\item Die \emph{Domänenschicht} enthält die Eignungsregeln und die Risikologik.
\item Die \emph{Datenschicht} enthält die Repositories für den Datenzugriff.
\end{itemize}
Die UI-Schicht folgt dem Android-Standardmuster Model-View-ViewModel (MVVM): Die Screens beobachten den Zustand ihrer ViewModels und melden Nutzeraktionen dorthin zurück. KI-Modell und Speicher liegen hinter Schnittstellen und bleiben so austauschbar.

\subsection{Bewertungskriterien}

Alle Komponenten werden nach denselben fünf Kriterien bewertet:
\begin{itemize}
\item \textbf{Datenschutz:} Wie viele Gesundheitsdaten verlassen das Gerät?
\item \textbf{Offline-Fähigkeit:} Funktioniert der Anwendungsfall ohne Netzverbindung?
\item \textbf{Aufwand:} Wie groß sind Umsetzung und Einarbeitung im Rahmen des Moduls?
\item \textbf{Kosten:} Welche laufenden Ausgaben für Betrieb und Hosting entstehen?
\item \textbf{Eignung:} Deckt die Lösung den Anwendungsfall vollständig ab?
\end{itemize}
Scheidet eine Option an einem Kriterium eindeutig aus, werden die übrigen nicht weiter betrachtet.

\subsection{Evaluation und Entscheidung}
\label{sec:evaluation}

\paragraph{Benutzeroberfläche: Kotlin mit Jetpack Compose.} Das MVP unterstützt nur Android (\secref{zielgruppe-primaer}), plattformübergreifende Frameworks wie Flutter bringen deshalb keinen Vorteil. Compose ist zudem schlanker als klassische XML-Layouts.

\paragraph{Kameraaufnahme: CameraX.} Die System-Kamera erlaubt kein Overlay und damit keine geführte Aufnahme. Camera2 ist aufwendig und anfällig für gerätespezifische Fehler. CameraX liefert Vorschau, Aufnahme und Bildanalyse als fertige Bausteine~\cite{android-camerax}.

\paragraph{Qualitätsprüfung: eigene Metriken.} OpenCV wäre für wenige Funktionen zu groß, ein eigenes KI-Modell bräuchte gelabelte Trainingsdaten. Für das MVP genügen zwei einfache Messwerte: die Schärfe über die Varianz des Laplace-gefilterten Bildes~\cite{pechpacheco2000} und die Belichtung über die mittlere Helligkeit.

\paragraph{Eignungsprüfung: Regeln in der Domänenschicht.} Die Ausschlussregeln sind feste Bedingungen, für die weder eine Bibliothek noch ein Dienst nötig ist.

\paragraph{KI-Einstufung: LiteRT auf dem Gerät.} Die Einstufung kann auf dem Gerät, auf einem eigenen Server oder über einen externen Dienst laufen. \tabref{ki-optionen} vergleicht die drei Optionen.

\begin{table}[htbp]
\caption{Vergleich der Optionen für die KI-Einstufung}
\label{tab:ki-optionen}
\footnotesize
\renewcommand{\arraystretch}{1.15}
\begin{tabularx}{\textwidth}{@{}>{\RaggedRight\arraybackslash}p{2.3cm} >{\RaggedRight\arraybackslash}X >{\RaggedRight\arraybackslash}X >{\RaggedRight\arraybackslash}X@{}}
\toprule
\textbf{Kriterium} & \textbf{LiteRT auf dem Gerät} & \textbf{Eigener Server (z.\,B. FastAPI mit ONNX Runtime)} & \textbf{Externer Dienst} \\
\midrule
Datenschutz & Fotos bleiben auf dem Gerät & Fotos gehen an den eigenen Server & Fotos gehen an Dritte \\
Offline-Fähigkeit & ja & nein & nein \\
Aufwand & Modell konvertieren und einbinden & zusätzlich Dienst und Hosting betreiben & gering, aber uns ist kein passender, frei nutzbarer Dienst bekannt \\
Kosten & keine & Hosting & nutzungsabhängig \\
Eignung & nur kompakte Modelle & auch große Modelle & vom Anbieter abhängig \\
\bottomrule
\end{tabularx}
\end{table}

Ausschlaggebend sind Datenschutz und Offline-Fähigkeit: Die Fotos verlassen das Smartphone nicht, und es muss kein Server betrieben werden. Backend-as-a-Service-Plattformen wie Supabase scheiden aus, weil ihre Edge Functions für Bildanalyse zu stark begrenzt sind~\cite{supabase-limits}.

LiteRT ist die Laufzeitumgebung, die Google bis 2024 als TensorFlow Lite entwickelt hat~\cite{litert2024}. Darin läuft ein für Mobilgeräte entworfenes Modell wie MobileNet~\cite{howard2017mobilenets}, das wir auf HAM10000~\cite{tschandl2018ham10000} oder dem ISIC-Archiv~\cite{isic-archive} nachtrainieren (\secref{machbarkeit}). Die drei Risikostufen entstehen über Schwellenwerte auf der Modellausgabe.

Zwei Einschränkungen nehmen wir bewusst in Kauf. HAM10000 enthält Nahaufnahmen vom Hautarzt, keine Smartphone-Fotos~\cite{tschandl2018ham10000}, und dunkle Hauttypen sind in öffentlichen Datensätzen kaum vertreten~\cite{wen2022datasets}. Das stützt den Ausschluss dunklerer Hauttypen in der Eignungsprüfung.

\paragraph{Datenhaltung: Room und privater App-Speicher.} Supabase würde Synchronisation und Gerätewechsel ermöglichen, verlangt dafür aber Login und die Speicherung von Gesundheitsdaten auf einem Server. Beides geht über das MVP hinaus. Die App speichert deshalb alles lokal: Strukturierte Daten liegen in einer Room-Datenbank~\cite{android-room}, Fotos im privaten App-Speicher. Von der Android-Datensicherung sind beide ausgeschlossen~\cite{android-autobackup}, damit keine Gesundheitsdaten ungewollt in die Cloud gelangen.

\paragraph{Körperübersicht: Compose Canvas.} Eine 3D-Bibliothek wäre für das MVP zu aufwendig, eine SVG-Grafik in einer WebView bräuchte eine Brücke zwischen Web- und App-Code. Auf dem Canvas lässt sich ein Körperumriss für Vorder- und Rückseite ohne zusätzliche Abhängigkeit zeichnen. Als Rückfallebene dient eine Auswahlliste der Körperregionen (\secref{machbarkeit}).

\paragraph{Erinnerungen: WorkManager.} Push-Nachrichten über Firebase Cloud Messaging brauchen einen Server. AlarmManager ist unnötig genau und verlangt ab Android 12 eine widerrufbare Sonderberechtigung~\cite{android-alarms}, obwohl die Intervalle Wochen bis Monate betragen. WorkManager plant Aufgaben auch nach einem Neustart des Geräts wieder ein~\cite{android-workmanager}.

\subsection{Gewählter Stack}

\tabref{stack} fasst die Entscheidungen zusammen.

\begin{table}[htbp]
\caption{Gewählter Technologiestack}
\label{tab:stack}
\footnotesize
\renewcommand{\arraystretch}{1.15}
\begin{tabularx}{\textwidth}{@{}l l X@{}}
\toprule
\textbf{Komponente} & \textbf{Anwendungsfälle} & \textbf{Technologie} \\
\midrule
Benutzeroberfläche & alle & Kotlin, Jetpack Compose, MVVM \\
Kameraaufnahme & UC-03 & CameraX \\
Qualitätsprüfung & UC-04 & eigene Metriken \\
Eignungsprüfung & UC-02, UC-07 & Regeln in der Domänenschicht \\
KI-Einstufung & UC-05, UC-06 & LiteRT, nachtrainiertes Modell \\
Datenhaltung & UC-01, UC-08, UC-09, UC-12 & Room, privater App-Speicher \\
Körperübersicht & UC-08, UC-09 & Compose Canvas \\
Erinnerungen & UC-10, UC-11 & WorkManager \\
\bottomrule
\end{tabularx}
\end{table}

Die optionalen Funktionalitäten F9 und F10 aus \secref{nicht-mvp} ändern an dieser Architektur nichts Grundlegendes. Die Arztsuche kann eine installierte Karten-App über einen Intent öffnen, und die UV-Warnung benötigt nur den Abruf eines öffentlichen UV-Index-Dienstes. Ein eigenes Backend ist für keinen der beiden nötig.
